% Generated draft: requires mathematical and source review.
\documentclass[11pt]{article}
\usepackage[margin=25mm]{geometry}
\usepackage{fontspec}
\setmainfont{Latin Modern Roman}
\usepackage{amsmath,amssymb,mathtools,unicode-math}
\setmathfont{Latin Modern Math}
\usepackage{enumitem,needspace,xurl,hyperref}
\setlength{\parindent}{0pt}
\setlength{\parskip}{5pt}
\setlength{\emergencystretch}{3em}
\newcommand{\N}{\mathbb{N}}
\newcommand{\Z}{\mathbb{Z}}
\newcommand{\Q}{\mathbb{Q}}
\newcommand{\R}{\mathbb{R}}
\newcommand{\C}{\mathbb{C}}
\newcommand{\F}{\mathbb{F}}
\newcommand{\A}{\mathbb{A}}
\newcommand{\PP}{\mathbb{P}}
\newcommand{\E}{\mathbb{E}}
\newcommand{\eps}{\varepsilon}
\newcommand{\dd}{\,\mathrm{d}}
\begin{document}
\section*{UnsolvedMath ID: 30004874}
\textbf{Source problem code:} catalogue entry
\subsection{Definitions and mathematical statement}
Let \(G\) be a compact connected simple Lie group, and let \(m_G\) be its normalized Haar probability measure. Let \(\mu\) be a Borel probability measure on \(G\). Write \(S=\operatorname{supp}(\mu)\) for its closed support. The measure is symmetric if \(\mu(A)=\mu(A^{-1})\) for every Borel set \(A\subseteq G\).

Define the convolution operator on \(L^2(G,m_G)\) by
\[
(\lambda_G(\mu)f)(x)=\int_G f(g^{-1}x)\,\mathrm{d}\mu(g),
\]
and define its invariant mean-zero subspace by
\[
L^2_0(G)=\left\{f\in L^2(G,m_G):\int_G f\,\mathrm{d}m_G=0\right\}.
\]
As in Definition 3 of the supplied excerpt, \(\mu\) has a spectral gap if
\[
\left\|\left.\lambda_G(\mu)\right|_{L^2_0(G)}\right\|_{\mathrm{op}}<1.
\]
Here the operator norm is the supremum of \(\|\lambda_G(\mu)f\|_2\) over all \(f\in L^2_0(G)\) with \(\|f\|_2=1\).

The excerpt does not explicitly define ``non-degenerate.'' Its explanation identifies containment of the support in a proper closed subgroup as the obstruction. Using that explanation, interpret non-degeneracy as
\[
\overline{\langle S\rangle}=G,
\]
where \(\langle S\rangle\) is the subgroup generated by \(S\). This interpretation is explicit but requires confirmation against the complete report: the accompanying lemma for general compact groups invokes products of two support elements, and the supplied literature assessment also uses the term ``adapted.''

\textbf{Spectral Gap Conjecture (supplied Conjecture 4).} For every compact connected simple Lie group \(G\) and every symmetric Borel probability measure \(\mu\) on \(G\), is it true that
\[
\left\|\left.\lambda_G(\mu)\right|_{L^2_0(G)}\right\|_{\mathrm{op}}<1
\quad\Longleftrightarrow\quad
\overline{\langle\operatorname{supp}(\mu)\rangle}=G?
\]
No finite-support, absolute-continuity, or algebraicity hypothesis is supplied. The convention for ``simple Lie group,'' including whether a nontrivial finite center is allowed, is not specified in the record.

The surrounding excerpt describes spectral gap as a uniform bound over nontrivial irreducible representations and states that it yields exponential equidistribution of convolution powers: for some \(c>0\) and a Sobolev norm \(S(f)\) of sufficiently high order,
\[
\int_G f\,\mathrm{d}\mu^{*n}
=\int_G f\,\mathrm{d}m_G+O\!\left(S(f)e^{-cn}\right)
\qquad(f\in C^\infty(G)).
\]
The precise Sobolev order and implicit-constant dependencies are not supplied. These are contextual source claims, not additional hypotheses. The extracted exponential-mixing formula has an unclear limiting term and is not adopted as a mathematical assertion here.
\subsection{Short English statement}
The supplied conjecture asks whether every symmetric probability measure whose support generates a dense subgroup of a compact connected simple Lie group has a convolution operator with norm strictly below one on mean-zero square-integrable functions, and conversely. The interpretation of ``non-degenerate'' and the convention for simplicity require source confirmation. The record labels the conjecture open; this draft does not independently verify that status.
\subsection{Sources}
\begin{itemize}
\item Supplied local UnsolvedMath registry record: permanent ID \(30004874\), exact source problem code \texttt{catalogue entry}, title ``Spectral-Gap Conjecture for Compact Simple Lie Groups,'' in the Oberwolfach Reports Open Problems collection. The excerpt identifies Definition 3 and Conjecture 4. No individual proposer is supplied; the recorded proposed year is 2021.
\item Supplied report citation: \emph{Arbeitsgemeinschaft: Thin Groups and Super-approximation} (2022), Oberwolfach Reports, pp.\ 2721--2790, DOI \texttt{10.4171/owr/2021/50}. \url{https://doi.org/10.4171/owr/2021/50}. The record attributes the problem to workshop open-problem discussions documented in this report.
\item Supplied literature reference, labeled ``Benoist--de Saxce'': \url{https://doi.org/10.1007/s00222-015-0636-2}.
\item Supplied literature reference, labeled ``preprint'': \url{https://arxiv.org/abs/1405.1808}.
\item Supplied literature reference, labeled ``OWR report'': \url{https://doi.org/10.4171/owr/2021/50}.
\item Literature assessment dated exactly \texttt{2026-08-22}: supplied dated curation, not independently verified. It labels the problem open, attributes to Benoist and de Saxce an equivalence between spectral gap and the almost-Diophantine property, and reports that whether every adapted, hence non-degenerate, measure is almost Diophantine is unknown. The supplied curation therefore describes the extracted conjecture as unresolved and explicitly requests independent verification of time-sensitive status and cited claims. Definitions and precise hypotheses of the almost-Diophantine equivalence are not included in the record.
\item The registry labels its statement \texttt{corrected\_verified} and reports that the official report was checked directly and extraction artifacts removed. This is supplied provenance, not a claim of source inspection or mathematical verification by this draft.
\end{itemize}
\end{document}
